% page 65 du fac-simile (image p-064 du scan)
% bloc 165.1 x 233.3 mm ; marge gauche 36.9 mm
\pgc{15.240mm}{0.000mm}{
\l{\cel{55}57}
\l{}
\l{\sou{baraktar}. (\sou{netrans.})\cel{2}Luktar per esforcar por su desintrikar.-R.}
\l{}
\l{\sou{baratrio}. Fraudo da la navestro kontre armatori od asekuristi.}
\l{\cel{3}- DEFIRS.}
\l{}
\l{\sou{barbo}.- Pili qui garnisas la mentono, la vangi, la labio supere}
\l{\cel{3}di adultulo. (\sou{anke metaf.})\cel{2}- DEFIRS.}
\l{}
\l{\sou{barbakano}. Fenduro facita, por la fluo de la pluv-aquo,}
\l{\cel{3}en muro qua subtenas tero. - DEFIS.}
\l{}
\l{\sou{barbara}. Stranjera pri civilizeso. - DEFIRS.}
\l{}
\l{\sou{barbelo}. (\sou{zool.}) River-fisho, ek la familio \textquotedbl{}ciprinoidi\textquotedbl{}, di qua}
\l{\cel{3}la maxilo esas garnisita ye filamento simila a ta qua existas}
\l{\cel{3}cirkum la boko di preske omna insekti - e di qua la karno}
\l{\cel{3}nutro-prizesas. - DEFIRS.}
\l{}
\l{\sou{barbotar}.\cel{2}(\sou{netrans.})\cel{2}Agitar su en aquo, igante lu spricar,}
\l{\cel{3}trublante lu. - FIS.}
\l{}
\l{\sou{barbuliar}.\cel{2}(\sou{netrans.}) Parolar tro rapide e, konseque, neklare}
\l{\cel{3}e konfuze, nekompreneble. - FIS.}
\l{}
\l{\sou{bardo}. Poeto Keltika qua kompozis e recitis kansoni militala,}
\l{\cel{3}religiala. - DEFIS.}
\l{}
\l{\sou{bardano}. (\sou{bot.}) Planto ek la familio \textquotedbl{}kompozaji\textquotedbl{}, di qua la ra-}
\l{\cel{3}diko uzesis olim medicinale kom sudorifigivo, e la (larja)}
\l{\cel{3}folii kom medikamento kontre tinio. - L.\cel{1}\sou{lappa}. - FIS.}
\l{}
\l{\sou{barejo}. Stofo de lano lejera, texita nedense. - DEFIRS.}
\l{}
\l{\sou{barelo}. Granda vazo de ligno, cilindra, inflita en la mezo,}
\l{\cel{3}konsistante ye duvi asemblita e cirkondita da ringi fera o}
\l{\cel{3}ligna, klozita per du fundi ligna, uzma por transportar}
\l{\cel{3}liquidi ed ula vari altra. - DEFIS.}
\l{}
\l{\sou{bario}. (\sou{kemio}). Korpo simpla, metala, flavatre blanka e brilanta,}
\l{\cel{3}tre oxidigebla, qua deskompozas la aquo ye la temperaturo}
\l{\cel{3}normala. -\cel{1}\sou{simbolo kemiala}\cel{1}: Ba. - DEFIS.}
\l{}
\l{\sou{baricentro}.\cel{2}Punto di korpo per qua pasas konstante la rezultanto}
\l{\cel{3}di la gravito-forci qui efikas ad omna ek lua molekuli. - FIS.}
\l{}
\l{\sou{barikado}.\cel{2}Fortifikuro improvizita, ek pavi amasigita, arbori}
\l{\cel{3}abatita, veturi renversita, e c., por impedar la aceso di}
\l{\cel{3}strado. - DEFIRS.}
\l{}
\l{\sou{barito}. (\sou{kemio}) Oxido di bario, substanco alkalioza venenoza,}
\l{\cel{3}saporanta korodive, quarople plu densa kam aero. -\cel{1}\sou{simbolo}}
\l{\cel{3}\sou{kemiala}\cel{1}: BaO.}
}
